\documentclass[12pt,a4paper]{article}
\usepackage[T2A]{fontenc}
\usepackage[utf8]{inputenc}
\usepackage[russian]{babel}
\usepackage{amsmath,amssymb,mathtools}
\usepackage[margin=22mm]{geometry}
\usepackage{microtype}
\usepackage{enumitem}
\usepackage{hyperref}
\hypersetup{hidelinks}
\setlength{\parindent}{0pt}
\setlength{\parskip}{0.45em}
\allowdisplaybreaks
\title{Дифференциальные уравнения\\\large Подготовка к контрольной}
\author{}
\date{}
\begin{document}
\maketitle
\textbf{Темы:} однородные и приводимые к ним уравнения; линейные уравнения первого порядка; уравнение Бернулли; метод введения параметра; уравнения Лагранжа и Клеро.

\textbf{Источник.} А.\,Ф.~Филиппов, «Сборник задач по дифференциальным уравнениям», издание 2000 года, §§\,2, 4, 5, 8. Справочные указания дополнены пояснениями и примерами. Загруженный скан \texttt{FilippovDU.pdf} прочитать пока не удалось; соответствие изданий не проверено.

\section{Область определения и потерянные решения}
Решением уравнения $y'=f(x,y)$ на интервале $I$ называется дифференцируемая функция $y$, график которой лежит в области определения $f$ и для каждого $x\in I$ удовлетворяет уравнению.

Каждое деление на выражение, которое может обратиться в нуль, сужает рассматриваемую область. Нужно отдельно подставить случай нулевого делителя в исходное уравнение. Например, деление на $y$ может потерять решение $y\equiv0$. Деление на $x$ требует отдельно рассмотреть возможные решения на интервалах $x>0$ и $x<0$.

В уравнении $M(x,y)\,dx+N(x,y)\,dy=0$ могут рассматриваться и кривые, не являющиеся графиками функций $y(x)$, например вертикальная прямая $x=0$. При переходе к $y'=-M/N$ они могут исчезнуть.

\section{Разделяющиеся переменные}
Если
\[
y'=f(x)g(y),
\]
то на участках, где $g(y)\ne0$, получаем
\[
\frac{y'}{g(y)}=f(x),
\qquad
\int\frac{dy}{g(y)}=\int f(x)\,dx+C.
\]
Здесь используется правило цепочки: если $G'(y)=1/g(y)$, то $\frac d{dx}G(y(x))=y'/g(y)$.

До деления найдём корни $g(y_0)=0$: при определённости уравнения они дают постоянные решения $y\equiv y_0$. При отсутствии единственности могут существовать также склеенные решения; это проверяется отдельно.

\textbf{Пример.} В уравнении $y'=2xy$ при $y\ne0$ имеем
\[
\frac{dy}{y}=2x\,dx,\qquad
\ln|y|=x^2+C,\qquad y=C_1e^{x^2}.
\]
Отдельно $y=0$ тоже является решением и включается в итоговую формулу при $C_1=0$.

\section{Однородные уравнения}
\subsection{Как распознать}
В форме, используемой Филипповым,
\[
y'=F\!\left(\frac yx\right).
\]
Правая часть зависит от $x,y$ только через отношение $y/x$. Другая форма:
\[
M(x,y)\,dx+N(x,y)\,dy=0,
\]
где $M,N$ однородны одной степени: $M(\lambda x,\lambda y)=\lambda^kM(x,y)$, аналогично для $N$, при $\lambda>0$.

\subsection{Замена $y=tx$}
Идея состоит в том, чтобы сделать отношение $y/x$ новой функцией:
\[
y=xt(x),\qquad y'=t+xt'.
\]
Следовательно,
\[
t+xt'=F(t),\qquad xt'=F(t)-t.
\]
Там, где $x\ne0$ и $F(t)-t\ne0$, переменные разделяются:
\[
\frac{dt}{F(t)-t}=\frac{dx}{x}.
\]
Если $F(t_0)=t_0$, то $t\equiv t_0$ даёт решение $y=t_0x$; оно могло потеряться при делении.

\textbf{Пример Филиппова.} Решим $x\,dy=(x+y)\,dx$. При $x\ne0$:
\[
t+xt'=1+t,\qquad t'=\frac1x,\qquad
t=\ln|x|+C,\qquad y=x(\ln|x|+C).
\]
В исходной дифференциальной форме отдельно имеется вертикальная интегральная кривая $x=0$; она не является графиком функции $y(x)$.

\textbf{Пример с потерянной ветвью.} Для $y'=(y/x)^2$ получаем $xt'=t(t-1)$. При $t\ne0,1$:
\[
\int\frac{dt}{t(t-1)}=\int\frac{dx}{x},
\qquad
\ln\left|\frac{t-1}{t}\right|=\ln|x|+C.
\]
После возвращения к $y$ получается $y=x/(1-Cx)$. Случаи $t=0,1$ дают $y=0$ и $y=x$; второе решение уже входит в семейство при $C=0$.

\section{Приведение к однородному}
\subsection{Перенос начала координат}
Пусть правая часть уравнения имеет вид
\[
y'=F\!\left(
\frac{a_1x+b_1y+c_1}{a_2x+b_2y+c_2}\right).
\]
Если прямые $a_1x+b_1y+c_1=0$ и $a_2x+b_2y+c_2=0$ пересекаются в $(x_0,y_0)$, вводим
\[
X=x-x_0,\qquad Y=y-y_0,\qquad \frac{dY}{dX}=y'.
\]
Свободные члены исчезают:
\[
\frac{dY}{dX}
=F\!\left(\frac{a_1X+b_1Y}{a_2X+b_2Y}\right).
\]
Это однородное уравнение, к нему применяем $Y=tX$.

\textbf{Если линейные части пропорциональны}, при различных параллельных прямых переноса в общую точку пересечения нет. Выбираем общую линейную часть $z=ax+by$; тогда $z'=a+by'$, и после подстановки часто получаем разделяющиеся переменные.

\subsection{Степенная замена}
Иногда помогает $y=z^m$, где $m$ подбирают из требования однородности после замены. При этом
\[
y'=m z^{m-1}z'.
\]
Необходимо учитывать область определения вещественной степени и отдельно проверять исключённый случай $z=0$.

\section{Линейное уравнение первого порядка}
\subsection{Форма и однородное уравнение}
Пусть $a,b$ непрерывны на интервале $I$:
\[
y'+a(x)y=b(x).
\]
Искомая функция и её производная входят линейно. Соответствующее линейное однородное уравнение получается при $b=0$.

Пусть $A'(x)=a(x)$. Тогда решения однородного уравнения:
\[
y_{\text{одн}}=Ce^{-A(x)}.
\]
Действительно, $y'_{\text{одн}}=-aCe^{-A}$.

\subsection{Вариация постоянной}
Филиппов предлагает заменить постоянную функцией:
\[
y=C(x)e^{-A(x)},\qquad
y'=C'e^{-A}-aCe^{-A}.
\]
Поэтому $y'+ay=C'e^{-A}=b$, откуда
\[
C'=be^A,\qquad
\boxed{y=e^{-A(x)}
\left(C_0+\int b(x)e^{A(x)}\,dx\right)}.
\]
Внутри скобок взята одна первообразная. Эквивалентно, $(e^Ay)'=e^Ab$: множитель $e^A$ превращает левую часть в производную произведения.

\textbf{Пример Филиппова, №\,136 в издании 2000 года.}
\[
xy'-2y=2x^4,\qquad x\ne0.
\]
Однородное уравнение имеет решение $y=Cx^2$. Положим $y=C(x)x^2$. Тогда
\[
x(C'x^2+2Cx)-2Cx^2=2x^4,
\quad C'=2x,\quad C=x^2+C_0.
\]
Итак, $y=x^4+C_0x^2$. Подстановка даёт $xy'-2y=2x^4$. Через точку $x=0$ продолжение надо рассматривать по исходному уравнению отдельно.

\subsection{Когда менять местами $x$ и $y$}
Если в дифференциальной форме выражение линейно относительно $x$ и $dx$, может оказаться удобным считать $x$ функцией $y$. На участке, где обратная функция существует и $y'\ne0$,
\[
\frac{dx}{dy}=\frac1{dy/dx}.
\]
Постоянные решения и точки, исключённые при делении на $y'$, проверяются отдельно.

\section{Уравнение Бернулли}
Пусть
\[
y'+a(x)y=b(x)y^n,\qquad n\ne0,1.
\]
При $n=0$ или $n=1$ уравнение линейное. На области, где $y\ne0$ и все рассматриваемые степени определены, делим на $y^n$:
\[
y^{-n}y'+a(x)y^{1-n}=b(x).
\]
Делаем замену $z=y^{1-n}$. По правилу цепочки
\[
z'=(1-n)y^{-n}y',\qquad
\boxed{z'+(1-n)a(x)z=(1-n)b(x).}
\]
Получилось линейное уравнение относительно $z$.

\textbf{Нулевая ветвь.} При $n>0$, если $y^n$ определено в нуле, функция $y\equiv0$ удовлетворяет исходному уравнению. При $n<0$ ноль не входит в область определения. При нецелом $n$ область определения степени следует оговорить отдельно.

\textbf{Пример Филиппова, №\,151 в издании 2000 года.}
\[
y'+2y=e^xy^2.
\]
Решение $y=0$ проверяем до деления. Для $y\ne0$ положим $z=1/y$:
\[
-z'+2z=e^x,\qquad z'-2z=-e^x.
\]
Подстановка $z=C(x)e^{2x}$ даёт
\[
C'e^{2x}=-e^x,\quad C'=-e^{-x},\quad
C=e^{-x}+C_0.
\]
Таким образом,
\[
z=e^x+C_0e^{2x},\qquad
y=\frac1{e^x+C_0e^{2x}}.
\]
Ответ рассматривается на интервалах, где знаменатель не равен нулю; отдельно добавляется $y=0$.

\section{Уравнения, не разрешённые относительно производной}
Если $F(x,y,y')=0$ удаётся представить как $y=f(x,y')$, положим
\[
p=y',\qquad y=f(x,p).
\]
Вдоль решения дифференцируем:
\[
p\,dx=f_x(x,p)\,dx+f_p(x,p)\,dp.
\]
Если на рассматриваемом участке можно выразить $x=x(p)$, то
\[
x=x(p),\qquad y=f(x(p),p)
\]
является кандидатом на параметрическое решение. При $x'(p)\ne0$ проверяем
\[
\frac{dy}{dx}=\frac{dy/dp}{dx/dp}=p.
\]
Первоначально $p$ — функция $x$, а параметром она становится локально на подходящих участках. Значения, потерянные при делениях, и постоянные $p$ разбираются отдельно. После дифференцирования ветви проверяются подстановкой в исходное уравнение.

\section{Уравнение Лагранжа}
\subsection{Вывод}
Уравнение Лагранжа имеет вид
\[
y=x\varphi(y')+\psi(y').
\]
Положив $p=y'$, получим $y=x\varphi(p)+\psi(p)$. Дифференцирование приводит к
\[
p\,dx=\varphi(p)\,dx+
\bigl(x\varphi'(p)+\psi'(p)\bigr)dp,
\]
или
\[
[p-\varphi(p)]\,dx
=\bigl(x\varphi'(p)+\psi'(p)\bigr)dp. \tag{1}
\]
Если $p=p_0$ постоянно, то из (1) следует $p_0=\varphi(p_0)$, и исходное уравнение даёт
\[
y=p_0x+\psi(p_0).
\]
Нельзя дописывать к этому выражению произвольную постоянную: исходное равенство уже определило свободный член.

Если $p-\varphi(p)\ne0$ и $x$ локально рассматривается как функция $p$, получаем линейное уравнение
\[
\frac{dx}{dp}
-\frac{\varphi'(p)}{p-\varphi(p)}x
=\frac{\psi'(p)}{p-\varphi(p)}.
\]
Решаем его и возвращаемся к $y=x\varphi(p)+\psi(p)$. Особые точки параметризации и возможные склейки необходимо анализировать отдельно.

\subsection{Пример}
Решим $y=2xy'-(y')^2$. С обозначением $p=y'$:
\[
y=2xp-p^2,\qquad -p\,dx=(2x-2p)\,dp.
\]
При $p\ne0$:
\[
\frac{dx}{dp}+\frac2p x=2,\qquad
\frac{d}{dp}(p^2x)=2p^2.
\]
Следовательно,
\[
x=\frac23p+\frac{C}{p^2},\qquad
y=2xp-p^2=\frac13p^2+\frac{2C}{p}.
\]
Проверка: $dy/dp=p\,dx/dp$. Там, где $dx/dp\ne0$, имеем $dy/dx=p$. Отдельно $p=0$ даёт решение $y=0$.

\section{Уравнение Клеро}
\subsection{Семейство прямых и огибающая}
Уравнение Клеро — частный случай Лагранжа:
\[
y=xy'+\psi(y').
\]
При $p=y'$ получаем $y=xp+\psi(p)$, а после дифференцирования
\[
\bigl(x+\psi'(p)\bigr)\,dp=0.
\]
Первая ветвь: $p=C$ постоянно, и
\[
\boxed{y=Cx+\psi(C).}
\]
Вторая ветвь задаётся параметрически:
\[
\boxed{x=-\psi'(p),\qquad
y=\psi(p)-p\psi'(p).}
\]
Это кандидат на огибающую семейства прямых. Нужно проверить, что соответствующая кривая задаёт график решения, а также проверить исходное уравнение. В вырожденных точках и при возможной склейке ветвей требуется отдельное исследование.

\subsection{Пример Филиппова, №\,287}
\[
y=xy'-(y')^2.
\]
Положив $p=y'$, получаем семейство $y=Cx-C^2$. Для второй ветви $x-2p=0$, поэтому
\[
p=\frac{x}{2},\qquad y=\frac{x^2}{4}.
\]
Проверка: $y'=x/2$ и $xy'-(y')^2=x^2/4$. Парабола касается прямой с параметром $C$ при $x=2C$: значения обеих функций равны $C^2$, а производные равны $C$.

\section{Памятка перед контрольной}
\begin{enumerate}[leftmargin=*]
\item Правая часть зависит от $y/x$: однородное, положить $y=tx$.
\item В дроби линейные выражения с константами: найти пересечение соответствующих прямых и перенести начало координат.
\item Вид $y'+a(x)y=b(x)$: решить однородное, затем варьировать постоянную.
\item Вид $y'+a(x)y=b(x)y^n$: проверить нулевую ветвь и положить $z=y^{1-n}$.
\item Вид $y=x\varphi(y')+\psi(y')$: Лагранж, положить $p=y'$; проверить $p=\varphi(p)$, затем решить линейное уравнение на $x(p)$.
\item Если $\varphi(p)=p$: Клеро; выписать семейство прямых и проверить огибающую.
\end{enumerate}
Для каждой задачи проверь область определения, нулевые делители и подстановку в исходное уравнение. Для параметрического ответа дополнительно проверь $dy/dx=p$.
\end{document}